\section{Mid-Term Examination --- Set B}
\label{sec:midterm_set_b}

\LPUPaperHeader{Mid-Term Examination (Objective MCQ)}{B}{90 Minutes}{30 Marks}{CO1 (Unit 1: Matrices), CO2 (Unit 2: ODEs \& Calculus), CO3 (Unit 3: Infinite Series)}

\begin{center}
\sffamily\textbf{\color{AlertRed}Instructions: All 30 questions are compulsory. Each question carries 1 mark. A penalty of 0.25 marks is deducted for each incorrect response.}
\end{center}

\vspace{3mm}

% ------------------------------------------------------------------------------
% 30 MCQS
% ------------------------------------------------------------------------------

\mcqitem{1}{If $A$ is a square matrix of order 3 such that $\det(A) = 4$, then $\det(2A)$ is:}{$8$}{$16$}{$32$}{$64$}{1}

\mcqitem{2}{The eigenvalues of a skew-symmetric matrix are always:}{Real}{Purely imaginary or zero}{Positive real numbers}{Complex with real part 1}{1}

\mcqitem{3}{The rank of the matrix $A = \begin{bmatrix} 1 & 2 & 3 \\ 2 & 4 & 6 \\ 3 & 6 & 9 \end{bmatrix}$ is:}{$3$}{$2$}{$1$}{$0$}{1}

\mcqitem{4}{If the system $x + y + z = 1, x + 2y + 4z = \alpha, x + 4y + 10z = \alpha^2$ is consistent, then $\alpha$ can be:}{$1$ or $2$}{$0$ or $1$}{$-1$ or $2$}{Any real number}{1}

\mcqitem{5}{If $\lambda$ is an eigenvalue of an invertible matrix $A$, then an eigenvalue of $\text{adj}(A)$ is:}{$\frac{|A|}{\lambda}$}{$\frac{\lambda}{|A|}$}{$\lambda |A|$}{$\frac{1}{\lambda}$}{1}

\mcqitem{6}{The trace of a $3\times 3$ matrix with eigenvalues $1, -2, 5$ is:}{$4$}{$-10$}{$6$}{$8$}{1}

\mcqitem{7}{The differential equation $(D^2 + 1)y = 0$ has the general solution:}{$y = c_1 e^x + c_2 e^{-x}$}{$y = c_1 \cos x + c_2 \sin x$}{$y = (c_1 + c_2 x)e^x$}{$y = c_1 \cosh x + c_2 \sinh x$}{1}

\mcqitem{8}{The value of $\lim_{x\to 0} \frac{\sin x - x}{x^3}$ is:}{$0$}{$\frac{1}{6}$}{$-\frac{1}{6}$}{$\frac{1}{3}$}{1}

\mcqitem{9}{The particular integral of $(D^2 - 1)y = x$ is:}{$-x$}{$x$}{$\frac{x^2}{2}$}{$e^x$}{1}

\mcqitem{10}{The general solution of $x y' + y = x^2$ is:}{$y = \frac{x^2}{3} + \frac{c}{x}$}{$y = x^3 + c$}{$y = \frac{x^3}{3} + c$}{$y = x + \frac{c}{x}$}{1}

\mcqitem{11}{The order and degree of the differential equation $\left[1 + \left(\frac{dy}{dx}\right)^2\right]^{3/2} = \frac{d^2 y}{dx^2}$ are:}{Order 2, Degree 3}{Order 2, Degree 2}{Order 1, Degree 3}{Order 3, Degree 2}{1}

\mcqitem{12}{The series $\sum_{n=1}^\infty \frac{n}{n^3 + 1}$ is:}{Convergent}{Divergent}{Oscillatory}{Conditionally Convergent}{1}

\mcqitem{13}{The series $\sum_{n=1}^\infty \frac{x^n}{n}$ converges for all $x$ in:}{$(-1, 1]$}{$[-1, 1)$}{$[-1, 1]$}{$(-1, 1)$}{1}

\mcqitem{14}{According to Cauchy's Root Test, $\sum u_n$ converges if $\lim_{n\to\infty} (u_n)^{1/n} = L$ where:}{$L < 1$}{$L > 1$}{$L = 1$}{$L \ge 1$}{1}

\mcqitem{15}{The series $\sum_{n=1}^\infty \frac{(-1)^{n-1}}{n^2}$ is:}{Conditionally Convergent}{Absolutely Convergent}{Divergent}{Unbounded}{1}

\mcqitem{16}{The derivative of $y = \tan^{-1}\left(\frac{2x}{1-x^2}\right)$ with respect to $x$ is:}{$\frac{1}{1+x^2}$}{$\frac{2}{1+x^2}$}{$\frac{2}{1-x^2}$}{$\frac{1}{1-x^2}$}{1}

\mcqitem{17}{The value of $\int_0^{\pi/2} \ln(\tan x) \, dx$ is:}{$0$}{$\frac{\pi}{2}$}{$\pi \ln 2$}{$-\frac{\pi}{2}\ln 2$}{1}

\mcqitem{18}{If $A = \begin{bmatrix} 1 & 0 \\ 0 & 1 \end{bmatrix}$, then the algebraic multiplicity of $\lambda = 1$ is:}{$1$}{$2$}{$0$}{$3$}{1}

\mcqitem{19}{If $A$ and $B$ are square matrices of the same order, then $\text{trace}(AB)$ is equal to:}{$\text{trace}(BA)$}{$\text{trace}(A)\text{trace}(B)$}{$\det(AB)$}{$0$}{1}

\mcqitem{20}{The value of $\int \sec^3 x \, dx$ is:}{$\frac{1}{2}[\sec x \tan x + \ln|\sec x + \tan x|] + C$}{$\sec x \tan x + C$}{$\frac{1}{3}\sec^3 x + C$}{$\tan^2 x + C$}{1}

\mcqitem{21}{The particular integral of $(D^2 + 2D + 1)y = e^{-x}$ is:}{$\frac{x^2}{2}e^{-x}$}{$x e^{-x}$}{$\frac{x}{2}e^{-x}$}{$\frac{x^3}{6}e^{-x}$}{1}

\mcqitem{22}{The test which is applied when D'Alembert's ratio test fails ($L=1$) is:}{Raabe's Test}{Leibniz's Test}{Integral Test}{Root Test}{1}

\mcqitem{23}{The series $\sum_{n=1}^\infty \frac{1}{n \ln n}$ ($n \ge 2$) is:}{Convergent}{Divergent}{Oscillatory}{Conditionally Convergent}{1}

\mcqitem{24}{The value of $\lim_{x\to 0} (1 + 2x)^{1/x}$ is:}{$1$}{$2$}{$e^2$}{$e$}{1}

\mcqitem{25}{A system of $m$ homogeneous linear equations in $n$ unknowns has a non-trivial solution if:}{$\text{rank}(A) = n$}{$\text{rank}(A) < n$}{$\text{rank}(A) > n$}{$\det(A) \neq 0$}{1}

\mcqitem{26}{If $A$ is idempotent ($A^2 = A$), its eigenvalues can only be:}{$0$ or $1$}{$1$ or $-1$}{$0$ or $-1$}{Any real number}{1}

\mcqitem{27}{If $A$ is nilpotent of index $k$ ($A^k = O$), all its eigenvalues are:}{$1$}{$-1$}{$0$}{$k$}{1}

\mcqitem{28}{The solution of the differential equation $\frac{dy}{dx} = \frac{y}{x}$ is:}{$y = c x$}{$y = c/x$}{$x^2 + y^2 = c$}{$y = x + c$}{1}

\mcqitem{29}{The Taylor series expansion of $\sin x$ about $x=0$ contains:}{Only even powers of $x$}{Only odd powers of $x$}{Both even and odd powers}{No powers of $x$}{1}

\mcqitem{30}{The geometric multiplicity of an eigenvalue $\lambda$ is always:}{$\le \text{Algebraic Multiplicity}$}{$> \text{Algebraic Multiplicity}$}{$= n$}{$= 0$}{1}

\vspace{5mm}
\hrule
\vspace{5mm}

% ------------------------------------------------------------------------------
% ANSWER KEY & EXPLANATIONS
% ------------------------------------------------------------------------------
\subsection*{Answer Key (Mid-Term Set B)}

\begin{center}
\begin{tabular}{|c|c||c|c||c|c||c|c||c|c||c|c|}
\hline
\textbf{Q} & \textbf{Ans} & \textbf{Q} & \textbf{Ans} & \textbf{Q} & \textbf{Ans} & \textbf{Q} & \textbf{Ans} & \textbf{Q} & \textbf{Ans} & \textbf{Q} & \textbf{Ans} \\
\hline
1 & \textbf{C} & 6 & \textbf{A} & 11 & \textbf{B} & 16 & \textbf{B} & 21 & \textbf{A} & 26 & \textbf{A} \\
2 & \textbf{B} & 7 & \textbf{B} & 12 & \textbf{A} & 17 & \textbf{A} & 22 & \textbf{A} & 27 & \textbf{C} \\
3 & \textbf{C} & 8 & \textbf{C} & 13 & \textbf{B} & 18 & \textbf{B} & 23 & \textbf{B} & 28 & \textbf{A} \\
4 & \textbf{A} & 9 & \textbf{A} & 14 & \textbf{A} & 19 & \textbf{A} & 24 & \textbf{C} & 29 & \textbf{B} \\
5 & \textbf{A} & 10 & \textbf{A} & 15 & \textbf{B} & 20 & \textbf{A} & 25 & \textbf{B} & 30 & \textbf{A} \\
\hline
\end{tabular}
\end{center}

\vspace{3mm}
\subsection*{Selected Detailed Mathematical Explanations}

\begin{solutionbox}[Explanations for Critical Questions]
\textbf{Q.1 (Ans: C)} For $n\times n$ matrix $A$, $\det(kA) = k^n \det(A)$. Here $n=3, k=2$: $\det(2A) = 2^3 \det(A) = 8(4) = 32$.\\[2mm]
\textbf{Q.3 (Ans: C)} Rows 2 and 3 are scalar multiples of Row 1 ($R_2 = 2R_1, R_3 = 3R_1$). Hence only 1 linearly independent row $\implies \text{rank} = 1$.\\[2mm]
\textbf{Q.5 (Ans: A)} Since $A \cdot \text{adj}(A) = |A|I \implies \text{adj}(A) = |A|A^{-1}$. If eigenvalue of $A$ is $\lambda$, eigenvalue of $A^{-1}$ is $1/\lambda$, so eigenvalue of $\text{adj}(A)$ is $|A|/\lambda$.\\[2mm]
\textbf{Q.8 (Ans: C)} $\lim_{x\to 0} \frac{(x - x^3/6 + \dots) - x}{x^3} = -\frac{1}{6}$.\\[2mm]
\textbf{Q.11 (Ans: B)} Squaring both sides gives $[1 + (y')^2]^3 = (y'')^2$. The highest derivative is $y''$ (order 2) and its power is 2 (degree 2).\\[2mm]
\textbf{Q.17 (Ans: A)} Using King's property $\int_0^a f(x)dx = \int_0^a f(a-x)dx$, $I = \int_0^{\pi/2} \ln(\cot x)dx = -\int_0^{\pi/2} \ln(\tan x)dx = -I \implies 2I = 0 \implies I = 0$.\\[2mm]
\textbf{Q.21 (Ans: A)} $\frac{1}{(D+1)^2}e^{-x} = \frac{x^2}{2!}e^{-x} = \frac{x^2}{2}e^{-x}$.\\[2mm]
\textbf{Q.24 (Ans: C)} $\lim_{x\to 0} (1+2x)^{1/x} = e^{\lim_{x\to 0} \frac{2x}{x}} = e^2$.
\end{solutionbox}

\newpage
